\section{Entry, transition, and telescoping}
\label{sec:telescoping}
%======================================================================

A source contributes at different times for different reasons.  Separating entry from transition shows which changes are genuine new data and which are internal transfers that cancel.

For a canonical source $m=cq$, define its first square-prefix index
\begin{equation}
 e(m):=\lfloor\sqrt m\rfloor.
 \label{eq:entry-index}
\end{equation}
This is the least $n$ such that $m<R_n^2$.  Define the smoothness-transition index
\begin{equation}
 h(m):=q-1.
 \label{eq:transition-index}
\end{equation}
This is the least $n$ such that $q\le R_n$.

The transport lifetime is
\begin{equation}
 L(m):=(h(m)-e(m))_+.
 \label{eq:lifetime}
\end{equation}
There are two cases:

\begin{enumerate}[leftmargin=1.7em,itemsep=0.3em]
\item If $q>\sqrt m$, equivalently $q>c$ and $Y>0$, then the point enters above the moving parabola and occupies transport for the discrete interval
\[
 e(m)\le n<h(m).
\]
\item If $q\le\sqrt m$, equivalently $q\le c$ and $Y\le0$, then $h(m)\le e(m)$ and the point enters already smooth.  Its transport interval is empty.
\end{enumerate}

Define the pointwise indicators
\begin{align}
 p_n(m)&:=\1_{\{e(m)\le n\}},
 \label{eq:p-indicator}\\
 u_n(m)&:=\1_{\{e(m)\le n<h(m)\}},
 \label{eq:u-indicator}\\
 a_n(m)&:=\1_{\{e(m)\le n,\ h(m)\le n\}}.
 \label{eq:a-indicator}
\end{align}
Then
\begin{equation}
 \boxed{p_n(m)=u_n(m)+a_n(m).}
 \label{eq:temporal-partition}
\end{equation}

\subsection{Why the crossing cancels}

Let $E_n^A$ be the signed mass entering the new vertical strip already below the moving parabola, let $E_n^U$ be the signed mass entering above it, and let $C_n$ be the signed mass crossed by the moving parabola between $R_{n-1}$ and $R_n$.  Then
\begin{align}
 \mathcal A(R_n)-\mathcal A(R_{n-1})
 &=E_n^A+C_n,
 \label{eq:Delta-A}\\
 \mathcal U(R_n)-\mathcal U(R_{n-1})
 &=E_n^U-C_n.
 \label{eq:Delta-U}
\end{align}
Therefore
\begin{equation}
 \boxed{
 S_n-S_{n-1}
 =E_n^A+E_n^U
 =\sum_{m\in B_n}\mu(m).
 }
 \label{eq:Delta-S}
\end{equation}
The entire crossing mass cancels.

\begin{figure}[htbp]
\centering
\begin{tikzpicture}[>=Latex,node distance=15mm]
  \node[draw,rounded corners,fill=gray!8,minimum width=30mm,minimum height=12mm,align=center] (out) {outside prefix\\$R<\sqrt{cq}$};
  \node[draw,rounded corners,fill=red!7,minimum width=32mm,minimum height=12mm,align=center,right=of out] (trans) {transport\\$\sqrt{cq}<R<q$};
  \node[draw,rounded corners,fill=green!8,minimum width=30mm,minimum height=12mm,align=center,right=of trans] (smooth) {smooth\\$q\le R$};
  \draw[->,very thick,navy] (out) -- node[above,align=center] {vertical entry\\$R=\sqrt{cq}$} (trans);
  \draw[->,very thick,teal] (trans) -- node[above,align=center] {parabolic crossing\\$R=q$} (smooth);
  \draw[decorate,decoration={brace,amplitude=6pt,mirror},thick] ($(trans.south west)+(0,-6mm)$) -- ($(trans.south east)+(0,-6mm)$) node[midway,below=7mm] {$\tau(c,q)=q-\sqrt{cq}$};
\end{tikzpicture}
\caption{The exact state path of a positive-height canonical point.  The second transition only reallocates signed mass between $-T$ and $A$, so it cancels from the residual increment.}
\label{fig:state-path}
\end{figure}

This is the telescopic nature of the squared coordinate: the point's imaginary height determines when the moving parabola reaches it, while the crossing itself contributes zero to the increment of the total prefix mass.

%======================================================================
